\documentclass[a4paper]{article}
\usepackage[T2A]{fontenc}
\usepackage[utf8]{inputenc}
\usepackage[russian]{babel}
\usepackage[margin=1.5cm]{geometry}
\usepackage{indentfirst}
\usepackage{amsmath}
\usepackage{amssymb}
\usepackage{amsthm}
\usepackage[hidelinks]{hyperref}


\theoremstyle{definition}
\newtheorem{definition}{Определение}[section]

\newtheorem{theorem}{Теорема}[section]

\title{Решение тестовых заданий Сбербанка}
\author{Берлет Максим}
\date{}

\begin{document}

\maketitle
\pagebreak

\tableofcontents
\pagebreak

\section{Задание 1}
\subsection{Условие}
Consider the simplest fair lottery: you make a bet V and with probability $\frac{1}{2}$ win (receive +V ) or lose (receive -V ). 
Consider following strategy $\{V_t\}_{t=1}^{\infty}$:
\[ 
V_t = \begin{cases}
    2^{t-1}, \ \text{if t = 1 or you lost all tries before t.} \\
    0, \ \text{otherwise.}
    \end{cases}
\]
In other words, you start with $V_1$ = 1 and double your bet until your first win, say at turn $\tau$, and don’t play
anymore after. Denote your profit (or debt) after t-th turn as $\pi_t$ with initial condition $\pi_0$ = 0. \\
1) Prove that $\tau$ < $\infty$ a.s.\\
2) Find $\mathbb{E}[\pi_{\tau}]$. \\
3) Find $\mathbb{E}[\pi_t]$. \\
4) Compare results in 2) and 3). Explain this in terms of Doob’s optional sampling theorem.

\subsection{Пункт 1}
    Требуется доказать, что рано или поздно найдется такой ход $\tau$, при котором произойдет выигрыш и игра закончится. 
    То есть, что вероятность такого события равна 1, при количестве ходов стремящемся к бесконечности.

    \begin{definition}
        Последовательность случайных величин $\xi_n$ сходится почти наверное (a.s.) к $\xi$ $\Leftrightarrow$ \\
        $P(\{\omega \in \Omega: \lim_{n \rightarrow \infty} \xi_n(\omega) = \xi(\omega) \}) = 1$  
    \end{definition}

    В данном случае $\tau$ является случайной величиной, которая зависит от исхода игры $\omega$. Покажем от обратного, что $\tau < \infty$.
    Пусть событие $A$ - игра никогда не закончится (то есть $\tau = \infty$), в таком случае $P(A) = 0$. За $L_t$ обозначим событие - проигрыш на ходе t.
    По условию $P(L_t) = \frac{1}{2}$, так как результаты каждого хода независимы, то вероятность проиграть n раз подряд равна $P = (\frac{1}{2})^n$. Иными словами
    вероятность, что выигрыш не наступит за n шагов $P(\tau > n) = (\frac{1}{2})^n$. При $n \rightarrow \infty$ получаем, что $P(\tau > \infty) = 0$. 
    Это значит, что вероятность того, что игра будет длиться вечно равна нулю, а значит $\tau < \infty$ почти наверное.     

\subsection{Пункт 2}
    В данном пункте требуется вычислить математическое ожидание прибыли в момент выигрыша $\tau$. 
    Рассмотрим конечные суммы $\pi_\tau$ при $\tau = k$. Мы можем разложить следующим образом $\pi_\tau = -V_{\text{проигрыша}} + V_{\text{выигрыша}} 
    = -(V_1 + V_2 + ... + V_{k-1}) + V_k = -(1 + 2 + ... + 2^{k-2}) + 2^{k-1}$, заметим, что первая скобка это в точности сумма геометрической прогрессии.
    $\pi_\tau = - (2^{k-1} - 1) + 2^{k-1} = 1$, таким образом получается, что $\pi_\tau = 1 = const$, т.е. прибыль не зависит от количества ходов и является константой, а математическое ожидание от константы равно ей же самой.
    $\mathbb{E}[\pi_\tau] = \mathbb{E}[1] = 1$.

\subsection{Пункт 3}
    В данном пункте требуется вычислить математическое ожидание кошелька не в момент победы, а в любой фиксированный ход.
    
    \begin{definition}
        Математическим ожиданием случайной величины $\xi$ называется $\mathbb{E}[\xi] = \int\limits_\Omega \xi(\omega) dP(\omega)$.
    \end{definition}

    Для нашего случая эквивалентно $\mathbb{E}[\xi] = \sum_{k=1}^{n} x_i * p_i$.
    Вычислим математическое ожидание для t = 1: \\ 
    $\mathbb{E}[\pi_1] = -1 * \frac{1}{2} + 1 * \frac{1}{2} = 0$. Вычислим для t = 2: $\mathbb{E}[\pi_2] = -2 * \frac{1}{2} + 2 * \frac{1}{2} + \mathbb{E}[\pi_1] = 0$. \\
    Выведем общую формулу: $\mathbb{E}[\pi_t] = -V_t * \frac{1}{2} + V_t * \frac{1}{2} + \mathbb{E}[\pi_{t-1}]$, а так как $\mathbb{E}[\pi_t] = \mathbb{E}[\pi_{t-1}] = ... = \mathbb{E}[\pi_1] = 0$.
    Это ожидаем, ведь игра по условию честная, то матожидание выигрыша на каждом отдельном ходу равно 0.

\subsection{Пункт 4}
    Сравним результаты из пункта 2 и пункта 3: $\mathbb{E}[\pi_\tau] = 1$, а $\mathbb{E}[\pi_t] = 0$.
    
    \begin{definition}
        Последовательность случайных величин $\{X_n\}_{n \in \mathbb{N}}$ - называется мартингалом, если: \\
        1. $\mathbb{E}[X_{n+1}| X_1, ..., X_n] = X_n, \ n \in \mathbb{N}$ \\
        2. $\mathbb{E}[|X_n|] < \infty, \ n \in \mathbb{N}$
    \end{definition}

    По определению наша последовательность случайных величин $\pi_t$ является мартингалом.

    \begin{theorem}[Теорема Дуба]
        Пусть $\{X_t\}_{t \in \mathbb{N} \cup \{0, \infty\}}$ - мартингал, а $\tau$ - время остановки, и предположим, что выполняется одно из следующих условий: \\
        1. Время остановки $\tau$ почти наверняка ограничено, то есть существует константа $c \in \mathbb{N}: \ \tau \le c$. \\
        2. Время остановки $\tau$ имеет конечное матожидание, а условные математические ожидания абсолютного значения приращений мартингала почти наверняка ограничены. \\
        3. Существует константа $c$ такая, что $|X_{min\{t, \tau\}}| < c$ почти наверняка для любого $t \in \mathbb{N}_0$ \\

        Тогда $\mathbb{E}[X_\tau] = \mathbb{E}[X_0]$.
    \end{theorem}

    Хотя последовательность $\pi_t$ и является мартингалом, однако она не удовлетворяет требованию на ограниченность случайной величины на любом ходу.
    Наш долг по модулю может быть неограниченно большим, а также нарушается требование 2 на ограниченность приращений, поэтому матожидание в момент остановки не равно матожиданию на первом ходу.

\section{Задание 2}
\subsection{Условие}
    Дана выборка из неизвестного распределения $\{x_i\}_{i=1}^{n}$. Упорядочим выборку $x_{(1)} \le x_{(2)} \le ... \le x_{(n)}$ и тогда мы имеем порядковые статистики $\{x_{(i)}\}_{i=1}^{n}$.
    Рассмотрим класс L-оценок: \\
    $L(x) = \sum\limits_{i=1}^{n} l_ix_{(i)}$, \\
    
    Оценка полностью определяется множеством параметров $\{l_i\}_{i=1}^{n}$. Замечание: некоторые $l_i$ могут быть равны нулю или быть функцией от $i$.\\
    Для следующих распределений предложите L-оценки параметров: \\
    (a) $\mathcal{U}[a, b]$, \\
    (b) $\mathcal{N}(\mu, \sigma)$, \\
    (c) $Pois(\lambda)$.

\subsection{Пункт a.}
    В данном пункте требуется найти L-оценку для параметров равномерного распределения $[a, b]$. 
    Так как фактически данные параметры задают границы, то логично за их оценки взять $a = x_{(1)},\ b = x_{(n)}$ соответственно.
    Проверим данную оценку на несмещенность.

    \begin{definition} 
        Оценка параметра $\theta$ называется несмещенной $\Leftrightarrow$ $\mathbb{E}[\hat{\theta}] = \theta$, где $\hat{\theta}$ - несмещенная оценка параметра $\theta$.
    \end{definition}

    Проверка на несмещенность: $\mathbb{E}[x_{(1)}] = \int\limits_{a}^{b}xp_{x_{(1)}}(x)dx$. \\
    Функция распределения для равномерного распределения: \\
    \[F_{X}(x) = 
        \begin{cases}
            0, \ x \le a \\
            \frac{x-a}{b-a}, \ a \le x < b, \\
            1, x \ge b
        \end{cases} = p(X < x)
    \] \\
    
    Для вычисления матожидания нам нужно найти функцию плотности для $x_{(1)}$: 
    $P(x_{(1)} > x) = P(x_1 > x, x_2 > x, ..., x_n > x) = P(X > x)^n$, так как наблюдения $x_i$ независимы.
    $P(X > x) = 1 - P(X < x) = 1 - \frac{x - a}{b - a} = \frac{b - x}{b - a}$, поэтому $P(x_{(1)} > x) = (\frac{b - x}{b - a})^n$. \\
    $F_{x_{(1)}}(x) = 1 - P(x_{(1)} > x) = 1 - (\frac{b - x}{b - a})^n$. Для вычисления плотности продифференцируем. \\
    $p_{x_{(1)}}(x) = n \frac{(b - x)^{n-1}}{(b - a)^n}$. 
    Вычислим матожидание: $\mathbb{E}[x_{(1)}] = \int\limits_{a}^{b}xn \frac{(b - x)^{n-1}}{(b - a)^n}dx = 
    \begin{bmatrix}
        u = x & du = dx \\
        dv = (b-x)^{n-1} & v = -\frac{(b-x)^n}{n} \\
    \end{bmatrix} = \frac{n}{(b-a)^n}\frac{-x(b-x)^n}{n}\Big|_{a}^{b} + \int\limits_a^b \frac{(b - x)^n}{n}dx = 0 + a + \frac{1}{n}\frac{b}{(b - a)^n}\frac{(b - x)^{n+1}}{n+1}(-1)\Big|_a^b = a + \frac{b-a}{n+1} = \frac{na + b}{n + 1}$
    Таким образом мы проверили оценку на несмещенность и данная оценка оказалась смещенной на величину $\frac{b - a}{n + 1}$. Проделаем аналогичные шаги для проверки оценки параметра $\hat{b} = x_{(n)}$.
    
    $F_{x_{(n)}}(x) = P(x_{(n)} \le x) = P(x_1 \le x, x_2 \le x, ..., x_n \le x) = P(X \le x)^n = (\frac{x - a}{b - a})^n$. \\
    $p_{x_{(n)}}(x) = \frac{dF}{dx} = n\frac{(x - a)^{n-1}}{(b - a)^n}$ \\
    $\mathbb{E}[x_{(n)}] = \int\limits_a^b xn\frac{(x - a)^{n-1}}{(b - a)^n} = 
    \begin{bmatrix}
        u = x & v = \frac{(x-a)^n}{n} \\
        dv = (x-a)^{n-1} & du = dx \\
    \end{bmatrix} = \frac{n}{(b - a)^n}(x \frac{(x - a)^n}{n}\Big|_a^b - \int\limits_a^b \frac{(x - a)^n}{n}dx) = 
    b - \frac{n}{(b - a)^n}\frac{1}{n}\frac{(x - a)^{n + 1}}{n + 1}\Big|_a^b = b - \frac{b - a}{n + 1} = \frac{bn + a}{n + 1}$.
    Данная оценка параметра $\hat{b} = x_{(n)}$ также оказалась смещенной на значение $\frac{b - a}{n + 1}$.

    Найдем несмещенную оценку параметра $a$ в виде $\hat{a} = l_1x_{(1)} + l_2x_{(n)}$. $\mathbb{E}[\hat{a}] = l_1\mathbb{E}[x_{(1)}] + l_2\mathbb{E}[x_{(n)}] = 
    l_1\frac{na + b}{n + 1} + l_2\frac{nb + a}{n + 1} = \frac{(l_1n + l_2)a + (l_1 + l_2n)b}{n+1}$. Так несмещенность оценки означает $\mathbb{E}[\hat{a}] = a$, 
    получаем систему: \\
    $
    \begin{cases}
        \frac{l_1n + l_2}{n + 1} = 1 \\
        \frac{l_1 + l_2n}{n + 1} = 0
    \end{cases}
    \Rightarrow l_1 = -l_2n \\
    -l_2n^2 + l_2 = n + 1 \\
    l_2(1 - n^2) = n + 1 \\
    l_2(1 - n)(1 + n) = n + 1\\
    l_2 = \frac{1}{1 - n} \\
    l_1 = \frac{n}{n - 1}
    $
    Получили несмещенную оценку $\hat{a} = \frac{n}{n - 1}x_{(1)} + \frac{1}{1 - n}x_{(n)} = \frac{nx_{(1)} - x_{(n)}}{n - 1}$.

    Проделаем аналогичные шаги для оценки $b$. \\
    $\hat{b} = l_1x_{(1)} + l_2x_{(n)}; \ \mathbb{E}[\hat{b}] = l_1\frac{na + b}{n + 1} + l_2\frac{nb + a}{n + 1} = \frac{(l_1n + l_2)a + (l_1 + l_2n)b}{n+1} $. \\
    $
    \begin{cases}
        \frac{l_1n + l_2}{n + 1} = 0 \\
        \frac{l_1 + l_2n}{n + 1} = 1
    \end{cases} 
    \Rightarrow
    l_2 = -l_1n \\
    l_1 + n(-l_1n) = n + 1 \\ 
    l_1 - n^2l_1 = n + 1 \\
    l_1 = \frac{1}{1 - n} \\
    l_2 = \frac{n}{n - 1}
    $
    Получили несмещенную оценку $\hat{b} = \frac{1}{1 - n}x_{(1)} + \frac{n}{n - 1}x_{(n)} = \frac{nx_{(n)} - x_{(1)}}{n - 1}$. \\

    $\hat{a} = x_{(1)} - \frac{x_{(n)} - x_{(1)}}{n - 1}$ \\
    $\hat{b} = x_{(n)} + \frac{x_{(n)} - x_{(1)}}{n - 1}$

\subsection{Пункт b.}
    В данном пункте требуется найти L-оценки для параметров нормального распределения $\mathcal{N}(\mu, \ \sigma)$. Самая очевидная L-оценка на параметр $\mu$ 
    это наилучшая среди несмещенных оценок - выборочное среднее, где все $l_i = \frac{1}{n}$, то есть $\hat{\mu} = \frac{1}{n}\sum\limits_{i=1}^nx_{(i)}$, которую легко можно получить методом максимального правдоподобия.
    Кроме того, среднее можно оценить медианой и данная оценка будет более устойчивой к выбросам(экстремальным значениям). \\
    $
    \hat{\mu} =
    \begin{cases}
        x_{([\frac{n}{2}] + 1)}, \ n \text{ - нечетное} \\
        \frac{x_{(\frac{n}{2})} + x_{(\frac{n}{2} + 1)}}{2}, \ n \text{ - четное}
    \end{cases}
    $

    Вывод оценки $\hat{\mu}$ методом максимального правдоподобия: $L(x_1, ..., x_n|\mu, \sigma) = \mathcal{L}(\mu, \sigma) = \prod\limits_{k=1}^{n}p(x_k|\mu, \sigma) = \\
    = \prod\limits_{k=1}^{n} \frac{1}{\sqrt{2\pi\sigma^2}} e^{-\frac{(x_k - \mu)^2}{2\sigma^2}} \Rightarrow log\mathcal{L}(\mu, \sigma) = 
    \sum\limits_{k=1}^{n} log(\frac{1}{\sqrt{2\pi\sigma^2}}) + log(e^{-\frac{(x_k - \mu)^2}{2\sigma^2}}) \Rightarrow \frac{\partial log\mathcal{L}}{\partial \mu} = \sum\limits_{k=1}^{n}\frac{2(x_k - \mu)}{2\sigma^2} = 0 \Rightarrow \frac{1}{n}\sum\limits_{k=1}^{n} x_k = \mu$

    В качестве оценки $\sigma$ можно взять размах: $x_{(n)} - x_{(1)} = \sigma(z_{(n)} - z_{(1)}), \ \text{где } z_i \sim \mathbb{N}(0, 1) \Rightarrow \\
    \mathbb{E}[x_{(n)} - x_{(1)}] = \sigma \mathbb{E}[z_{(n)} - z_{(1)}]$, то есть оценка смещена на матожидание размаха в стандартном нормальном распределении, а значит несмещенная оценка будет выглядеть следующим образом: \\
    $\hat{\sigma} = \frac{x_{(n)} - x_{(1)}}{d_n}$, где $d_n = \mathbb{E}[z_{(n)} - z_{(1)}]$ зависит только от n.

\subsection{Пункт c.}
    В данном пункте требуется найти L-оценку параметра $\lambda$ Пуассоновского распределения $Pois(\lambda)$. Очевидная оценка для $\lambda$ является выборочное среднее, кроме того, такая оценка является ещё и лучшей среди несмещенных. 
    Получить её можно также с помощью метода максимального правдоподобия, покажем: $L(x_1, ..., x_n | \lambda) = \mathcal{L}(\lambda) = \\ 
    \prod\limits_{k=1}^n p(x = x_k| \lambda) = \prod\limits_{k=1}^{n} \frac{\lambda^{x_k}e^{-\lambda}}{x_k!} \Rightarrow log\mathcal{L}(\lambda) = \sum\limits_{k=1}^{n} log(\lambda^{x_{k}}e^{-\lambda}) - log(x_k!) \Rightarrow 
    \frac{\partial log\mathcal{L}}{\partial \lambda} = \sum\limits_{k=1}^n \frac{x_k}{\lambda} - 1 = 0 \Rightarrow \lambda = \frac{1}{n}\sum\limits_{k=1}^n x_k$ 
\section{Задание 3}
\subsection{Условие}
    Рассмотрите следующую опционную стратегию. Вам разрешено купить $N_1$ колл-опционов со страйком 90\% от ATM, продать 
    $N_2$ колл-опционов со страйком 110\% от ATM и купить или продать некоторое количество базового актива $N_3$. 
    Все опционы имеют одинаковый срок до погашения $T$.

    1) Используйте модель Garman–Kohlhagen для оценки всех вышеуказанных опционов на валютной паре USD/EUR с $T = 1$ год. 
    Объясните, как вы задаёте параметры модели и почему именно так.

    2) При фиксированном общем количестве $N = \sum_{i=1}^{3} N_i^*$, найдите оптимальные пропорции капитала $\left( \frac{N_1^*}{N}, \frac{N_2^*}{N}, \frac{N_3^*}{N} \right)$, 
    при которых $\Delta$ и $\Gamma$ стратегии сбрасываются в ноль.

    3) Предположим, вы управляете капиталом в размере USD 100,000 и полностью инвестируете его в оптимальную стратегию из пункта 2).
    Рассчитайте $\Theta$ портфеля и дайте интерпретацию результата. Какую доходность стратегии вы можете ожидать?

    4) Какие предположения должны выполняться, чтобы модель из пункта 1) работала? Что может пойти не так в реальной жизни? Как это можно исправить?
    
\subsection{Пункт 1}
    
\subsection{Пункт 2}

\subsection{Пункт 3}

\subsection{Пункт 4}

\end{document}